\newcommand{\DoNotLoadEpstopdf}{}
\documentclass[11pt,reqno]{amsart}
\usepackage[T1]{fontenc}
\usepackage{lmodern,amsmath,amssymb,amsthm,mathtools,mathrsfs,microtype}
\usepackage[margin=1.1in]{geometry}
\usepackage{xr-hyper}
\usepackage[hidelinks]{hyperref}
\externaldocument[supp-]{build/supplement}[supplement.pdf]
\hypersetup{pdftitle={Collision records, arithmetic, and raw local inversion in a triangular Lorentz gas},pdfauthor={Qian Qi}}
\setlength{\emergencystretch}{2em}
\allowdisplaybreaks[2]
\numberwithin{equation}{section}
\newtheorem{theorem}{Theorem}[section]
\newtheorem{lemma}[theorem]{Lemma}
\newtheorem{proposition}[theorem]{Proposition}
\newtheorem{corollary}[theorem]{Corollary}
\theoremstyle{definition}\newtheorem{definition}[theorem]{Definition}
\newtheorem{inputstatement}[theorem]{Input statement}
\theoremstyle{remark}\newtheorem{remark}[theorem]{Remark}
\newcommand{\R}{\mathbb R}\newcommand{\Z}{\mathbb Z}
\newcommand{\one}{\mathbf1}\newcommand{\dd}{\,\mathrm d}
\newcommand{\esssup}{\operatorname*{ess\,sup}}
\newcommand{\diver}{\operatorname{div}}\newcommand{\TV}{\operatorname{TV}}
\title[Collision records and raw local inversion]{Collision records, arithmetic, and raw local inversion in a triangular Lorentz gas}
\author{Qian Qi}
\date{October 14, 2026. A2-DYN, revision 107}
\subjclass[2020]{37D50, 37A50, 60F05, 28A75}
\keywords{Lorentz gas, exact collision records, positive paired traces, stationary endpoint roofs, physical residence, angular comparison, critical entrances}
\begin{document}
\raggedbottom
\begin{abstract}
We investigate the unchanged exact-label raw pointwise local limit
problem for actual returns in a triangular Lorentz gas. On the
original positive paired endpoint source, we prove a finite-count
persistence inequality that pays every physical sheet at a previously
fixed positive total residence, regardless of flatness, critical
endpoint structure, physical comparison reach or stationarity.
The corresponding essential-height constant is independent of the
fixed positive residence cutoff. We also compare actual stationary
critical arms at an adaptive, smaller physical reach and pay their
complete positive outgoing-trace shortage by the original uncut
optical volume. Combining the two mechanisms with the inherited
critical-arm comparison reduces the original pointwise error
to the original unmatched shrinking-residence positive source.
We formulate an explicit untruncated residence-moment sufficient
criterion for its decay without asserting that the Lorentz source
satisfies that missing criterion. The original pointwise theorem,
section, exact occupation and displacement labels, arithmetic
zero classes and scalar-first conditional-law order are unchanged.
We further prove a two-ended positive-envelope comparison on
nonstationary physical sheets. All losses of either endpoint trace
are paid by the original volume on a common roof window; only
unallocated comparison capacity remains in the positive remainder.
This comparison requires neither a stationary incoming roof nor a
positive lower bound for the paired trace.
\end{abstract}
\maketitle
\input{journal-problem-v93}

\subsection{Comparison at a stationary endpoint}
The previous endpoint-distance flows compare nearby trajectories of
the same exact event. Their paired coarea densities are
$c^\varepsilon/|a'|$ and $c^\varepsilon/|b'|$, where
$c^\varepsilon$ is the positive minimum of the weighted traces at
the two ends of a maximal regular segment. These formulas remove
the relative-position degeneracy but still require nonzero endpoint
speed along a useful comparison path.

The argument below uses intervals of actual angular trajectories
instead of a single comparison path. If all trajectories in a
comparison interval contain a one-sided roof translate of every
source trajectory, the integral of their positive trace controls the
source at its original roof. The ratio of the two integrated traces
and the physical overlap are the only numerical costs. Stationary
endpoint velocity and rapid local variation of the minimum do not
enter that ratio. Theorem~\ref{thm:v104-stationary} verifies it on
critical arms with an explicit level-width condition and the original uncut optical variation. Theorem~\ref{thm:v104-paid-arms}
also pays the entire comparison-trace shortage. Its whole-family
height bound includes all
relative positions, without a lower original residence cutoff.

The pointwise problem is still \eqref{eq:v76-target}. The new
theorems estimate specified positive source families at that topology,
not a replacement finite-norm problem. Theorem~\ref{thm:v104-recovery}
retains a nonnegative remainder with a precise comparison-trace
deficit. Its vanishing on the complete physical family remains the
last requirement in the resulting scalar reduction.

\subsection{Comparison without a stationary entrance}
The two endpoint traces play symmetric roles in the next comparison.
Their positive maximum supplies comparison capacity, while their
absolute difference is controlled by the original optical variation.
Theorem~\ref{thm:v107-envelope} proves the resulting essential-height
estimate on actual physical sheets with moving entrances and exits.
Lemma~\ref{lem:v107-window} localizes the entire comparison segment
near the target roof by interval inclusion alone. It replaces the
common-incoming-roof argument, rather than assuming its conclusion.
Theorem~\ref{thm:v107-allocation} then retains, with its original
weight, precisely the unallocated part of the microscopic source.
No physical coverage or microscopic height limit is inferred from
finite-count definability or from a model calculation.

\section{The inherited physical source in the required form}
Fix the radius, collision count and exact displacement and occupation
labels. On the disjoint original space of regular pairing sheets,
write $a(\theta)<u<b(\theta)$ for a maximal roof segment. On the regular set away from $g=0$, put
\[
 q=(\alpha_0,\alpha_m),\quad F=L,\quad
 g=\nabla_qF=R(-p_0,p_m),\quad z=|g|,\quad
 I_m=1+\sum_{i=1}^{m-1}q_0^{d_i}/c_i.
\]
Here $\omega=|\partial_{p_0}\alpha_m|^{-1}$ and
$\chi=|\det D_q(F,\arg g)|$. The physical angular density is
\[
 A(u,\theta)=\frac{\omega(u,\theta)}{4\pi c_*\chi(u,\theta)},
 \qquad f^\varepsilon=A W^\varepsilon,
 \qquad c^\varepsilon(\theta)
     =\min\{f^\varepsilon(a+,\theta),f^\varepsilon(b-,\theta)\}.
\]
The initial section visit is counted and the terminal section visit
is tested separately, as in \eqref{eq:v76-event}. Every angular
sheet carries this same exact event; its incident physical component
is part of its identity.

\begin{inputstatement}[Positive pairing and local concentration]
\label{input:v104-pair}
The original paired measure is
$c^\varepsilon(\theta)\,d\theta\,du$ on $a<u<b$.
For $0<\Delta$ let $\mathsf T_\Delta^\varepsilon$ be its roof
density restricted to $0<\ell=b-a<\Delta$, without an additional
selector. At every previously fixed $r,\Delta,\varepsilon>0$,
\begin{equation}\label{eq:v104-paired-average}
 \limsup_m\sup_{R,n,k}m^2
       \|A_r\mathsf T_\Delta^\varepsilon\|_\infty
 \le C\{\varepsilon^{1/16}+(1+\varepsilon^{-1})\Delta\},
 \qquad A_rf(t)=\frac1{2r}\int_{t-r}^{t+r}f(u)\,du.
\end{equation}
The leading constant is independent of the previously fixed widths.
For fixed $m,R,n,k$, all nonnegative angular restrictions use the
same coarea kernels outside a roof-null set. Countable selections
use common versions. The uniform order is $\sup_R\esssup_t$;
one universal null set for all radii is not asserted.
\end{inputstatement}
This is the positive barycentric source and complete paired-average
estimate of Supplement Theorem~\ref{supp-v103archive:thm:v102-pair-coarea}
and equation~\eqref{supp-v103archive:eq:v102-paired-average}.
The paired-average estimate follows there from the original weighted
source and the actual uncut inverse-gradient volume source. It is
not an assumed essential-height estimate for the paired source.
The same inherited analytic input gives the actual uncut source
\begin{equation}\label{eq:v104-uncut-input}
 v^0(t)\,dt=F_\#(I_m z^{-1}\one_\Pi\nu_R^*),\qquad
 \limsup_m\sup_{R,n,k}m^2\|A_\rho v^0\|_\infty\le C
\end{equation}
for every previously fixed $\rho>0$, with a width-independent leading
constant. It is finite at fixed count. On each regular segment,
\begin{equation}\label{eq:v104-weighted-variation-input}
 |\partial_u f^\varepsilon(u,\theta)|
 \le C(1+\varepsilon^{-1})A(u,\theta)I_m(u,\theta)/z(u,\theta).
\end{equation}
These are the actual uncut volume and weighted optical variation
statements in Supplement Input statement~\ref{supp-v103archive:input:v102-analytic}
and Input statement~\ref{supp-v103archive:input:v102-geometry}.
The latter records the maximal-segment map and one-sided trace
comparison; the underlying statements are Supplement
Theorem~\ref{supp-v85archive:thm:v85-segment-map} and
Theorem~\ref{supp-v100archive:thm:v98-paired-trace}.
The uncut bound is Supplement
Theorem~\ref{supp-v100archive:thm:v100-uncut}.
No unweighted pointwise height bound for
$v^0$ is assumed. The bound is used only on a previously fixed
positive window after the physical sum.

The Borel physical sheet summation retains multiplicity one; its
historical proof and both endpoint trace conventions are unchanged.

\begin{inputstatement}[Regular critical entrance and scalar reduction]
\label{input:v104-scalar}
At a regular nongrazing backward-critical extension $q_*$, the
incoming roof is constant and
\begin{equation}\label{eq:v104-critical-input}
 f^\varepsilon(a+,\theta)
 =\frac{N_*}{e_\perp(\theta)^{\mathsf T}K_*e_\perp(\theta)},
 \quad N_*=\frac{\omega(q_*)W^\varepsilon(q_*)}{4\pi c_*},
 \quad \lambda_\circ I\le K_*\le\Lambda_\circ I,
\end{equation}
where $\lambda_\circ=9\sqrt{19}/200$ and $\Lambda_\circ=47/6$.
The incoming trace can be positive. A singular stationary extension
is not included in this formula merely because its incoming roof
is critical.

Under the original outer schedule
\begin{equation}\label{eq:v104-schedule}
 \varepsilon(B)=A_0B^{-1/12},\qquad q_0^*=1/8,\qquad
 r_0=c_0\varepsilon(B)^4,\qquad\delta=6q_0^*r_0,
\end{equation}
the preceding positive remainder $\mathsf U_B^{\rm rem}$ is the
original short paired source on failure of the relative clock and
both endpoint-distance clocks. With
\[
 \mathcal H_M(T)=\limsup_m\sup_{R,n,k}m^2
       \esssup_{t\in\mathcal C_{m,R,n,k}(M)}T(t),
\]
one has
\begin{equation}\label{eq:v104-old-remainder}
 \mathcal E_M\le\mathcal H_M(\mathsf U_B^{\rm rem})
                                       +C_MB^{-1/384}.
\end{equation}
Along this same schedule, the original pointwise assertion is
equivalent to $\mathcal H_M(\mathsf U_B^{\rm rem})\to0$.
\end{inputstatement}
The trace formula is Supplement
Proposition~\ref{supp-v103archive:prop:v103-critical-trace}; the
scalar statement is Supplement Theorem~\ref{supp-v103archive:thm:v103-recovery}.
The earlier relative and endpoint-distance comparisons remain
available and unchanged. The new comparison is applied specifically
to their positive failure source below.

\input{core/250_positive_angular_interval_comparison}
\input{core/251_stationary_critical_arm_recovery}
\input{core/252_positive_deficit_and_original_endpoint}
\input{core/253_uniform_positive_residence_persistence}
\input{core/254_adaptive_reach_and_paid_trace_deficit}
\input{core/255_exact_microscopic_and_failure_partition}
\input{appendices/v106_logarithmic_microscopic_criterion}

\input{core/256_two_ended_envelope_and_moving_roof_comparison}
\input{core/257_fractional_capacity_and_original_microscopic_source}

\appendix
\input{appendices/v104_dependency_map}
\input{appendices/v105_dependency_map}
\input{appendices/v107_coverage_and_limit_order}
\input{journal-references-v100}
\end{document}

